\section{Conclusion}
\label{sec:conclusion-ra}

The amendment side of constitutional admission is type-conditioned on
a backward-refinement witness; the response side, for actions whose
effects are reversible by an inverse executor, can be conditioned on
a positive TTL witness and an optional rollback receipt that closes
the action chain. The composition theorem
\thm{ttl_bounded_amendment_chain_preserves_baseline} relates the
response side to the amendment side: a chain of TTL-bounded
amendments, each carrying a per-step refinement witness, preserves
the baseline constitution at every instant under the auto-reversion
semantics of the TTL window. The reduction
\thm{destructive_action_requires_bilateral_admission} specializes the
treaty intersection theorem of the prior substrate
\cite{programmableSovereigntyAnonymous} to the destructive
subclass of the action taxonomy, making human-in-the-loop oversight
of destructive enforcement a typed obligation rather than a policy
file. The endpoint enforcement engine grounds the substrate in four
reversible action variants with content-hash-gated inverse executors
and a four-receipt grammar; the TTL auto-expiry scheduler, the
missing inverse executors on the destructive subclass, and the
bilateral cosignature path remain substrate-level obligations the
runtime does not yet deploy. The headline composition theorem
discharges in the paper-local mechanization by a case split on the TTL
window with per-step witnesses on an explicit finite admission domain,
not by reflexivity; on that domain the response side carries the same
finite-domain refinement discipline as the amendment side of the prior
substrate. Promoting the mechanization into the substrate's proof root
is the next step.
